\documentclass[11pt]{article}
\usepackage[margin=1in]{geometry}
\usepackage{amsmath, amssymb, amsthm}
\usepackage{booktabs}
\usepackage{hyperref}

\newtheorem{theorem}{Theorem}
\newtheorem{lemma}[theorem]{Lemma}
\newtheorem{proposition}[theorem]{Proposition}

\title{Project Euler Problem 390:\\Triangles with Non-Rational Sides and Integral Area}
\author{}
\date{}

\begin{document}
\maketitle

\section{Problem Statement}

Consider the triangle with sides $\sqrt{5}$, $\sqrt{65}$, and $\sqrt{68}$.
It can be shown that this triangle has area~9.

Define $S(n)$ as the sum of the areas of all triangles with sides
$\sqrt{1+b^2}$, $\sqrt{1+c^2}$, and $\sqrt{b^2+c^2}$ (for positive
integers $b$ and~$c$) that have an integral area not exceeding~$n$.

Given: $S(50) = 9$.

Find: $S(10^{10})$.

\section{Area Formula}

\begin{proposition}
The area of the triangle with sides $\sqrt{1+b^2}$, $\sqrt{1+c^2}$,
$\sqrt{b^2+c^2}$ is
\[
  K = \frac{1}{2}\sqrt{b^2 c^2 + b^2 + c^2}.
\]
\end{proposition}

\begin{proof}
Embed the triangle in $\mathbb{R}^3$ with vertices
\[
  P_1 = (0,0,0), \quad P_2 = (1,b,0), \quad P_3 = (1,0,c).
\]
One verifies $|P_1P_2| = \sqrt{1+b^2}$, $|P_1P_3| = \sqrt{1+c^2}$,
$|P_2P_3| = \sqrt{b^2+c^2}$. The area equals half the magnitude of the
cross product:
\[
  \overrightarrow{P_1P_2} \times \overrightarrow{P_1P_3}
  = (1,b,0) \times (1,0,c) = (bc,\, -c,\, -b),
\]
so $K = \tfrac{1}{2}\sqrt{b^2c^2 + c^2 + b^2}$.
\end{proof}

\section{Integrality Conditions}

\begin{lemma}
$K \in \mathbb{Z}^+$ if and only if both $b$ and $c$ are even.
\end{lemma}

\begin{proof}
We need $b^2c^2 + b^2 + c^2 = (2K)^2$, a perfect square divisible by~4.
Examining $b^2c^2 + b^2 + c^2 \pmod{4}$:

\begin{center}
\begin{tabular}{cccc}
\toprule
$b \bmod 2$ & $c \bmod 2$ & $b^2c^2+b^2+c^2 \pmod{4}$ & Result \\
\midrule
odd & odd & 3 & never $\equiv 0$ \\
odd & even & 1 & root is odd \\
even & odd & 1 & root is odd \\
even & even & 0 & possible \\
\bottomrule
\end{tabular}
\end{center}
Only when both are even can the expression be divisible by~4.
\end{proof}

Writing $b = 2B$, $c = 2C$ with $B, C \ge 1$:
\begin{equation}\label{eq:area}
  K = \sqrt{4B^2C^2 + B^2 + C^2} \in \mathbb{Z}^+.
\end{equation}

\section{Reformulation as a Pell Equation}

\begin{proposition}
The condition~\eqref{eq:area} is equivalent to
\[
  (b^2+1)(c^2+1) = (2K)^2 + 1.
\]
\end{proposition}

\begin{proof}
$(b^2+1)(c^2+1) = b^2c^2 + b^2 + c^2 + 1 = 4K^2 + 1$.
\end{proof}

For fixed even $a = \min(b,c)$, define $D := a^2+1$. Writing
$b_{\text{other}} = at + v$ and $n := a \cdot b_{\text{other}} + t = 2K$,
we obtain the \textbf{generalized Pell equation}:
\begin{equation}\label{eq:pell}
  v^2 - D \cdot t^2 = -a^2,
  \qquad D = a^2+1.
\end{equation}

\section{Pell Group Structure}

The \emph{fundamental solution} of $x^2 - Dy^2 = 1$ is
\[
  (p, q) = (2a^2+1,\; 2a),
\]
since $(2a^2+1)^2 - (a^2+1)(2a)^2 = 1$.

\begin{lemma}
Since $D = a^2+1$, the equation $x^2 - Dy^2 = -1$ has the solution
$(x,y) = (a, 1)$.
\end{lemma}

This implies the full automorphism group of the Pell form is generated by
$(a, 1)$, with $(p, q) = (a,1)^2$.

\subsection{Orbit Representatives}

Solutions of~\eqref{eq:pell} fall into orbits under the action
$(v, t) \mapsto (vp + Dtq,\; vq + tp)$.

\begin{proposition}
Orbit representatives $(v_0, t_0)$ satisfy $t_0 \in \{2, 4, \ldots, a\}$
(even) with $Dt_0^2 - a^2$ a perfect square.
\end{proposition}

\subsection{Sign Doubling}

\begin{theorem}\label{thm:sign}
For a non-trivial orbit representative with $t_0 < a$, the orbits seeded by
$(+v_0, t_0)$ and $(-v_0, t_0)$ are \textbf{distinct}. For the trivial
representative $t_0 = a$, $v_0 = a^2$, they \textbf{coincide}.
\end{theorem}

\begin{proof}
The orbit of $(-v_0, t_0)$ coincides with that of $(v_0, t_0)$ if and only
if there exists $k \ge 1$ with
\[
  (-v_0, t_0) \cdot (p, q)^k = (v_0, t_0).
\]
For $k = 1$: the equations $-v_0 p + Dt_0 q = v_0$ and $-v_0 q + t_0 p = t_0$
yield $t_0(p-1) = v_0 q$, hence $t_0 \cdot 2a^2 = v_0 \cdot 2a$, giving
$v_0 = a \cdot t_0$. Substituting into $v_0^2 = Dt_0^2 - a^2$:
\[
  a^2 t_0^2 = (a^2+1)t_0^2 - a^2 \implies t_0^2 = a^2 \implies t_0 = a.
\]
For $k \ge 2$, the $t$-component grows strictly, precluding equality.
\end{proof}

\section{Algorithm}

\begin{enumerate}
  \item For each even $a$ from $2$ to $\lfloor\sqrt{2n}\rfloor$:
    \begin{enumerate}
      \item For each even $t_0$ from $2$ to $a$, check if
            $(a^2+1)t_0^2 - a^2$ is a perfect square $v_0^2$.
      \item If $t_0 < a$: iterate orbits from both $(v_0, t_0)$ and $(-v_0, t_0)$.
      \item If $t_0 = a$: iterate only from $(v_0, t_0) = (a^2, a)$.
      \item In each orbit, apply the Pell recurrence
            $(v, t) \mapsto (vp + Dtq,\; vq + tp)$ until the derived area exceeds~$n$.
      \item For each valid $(v, t)$ with $v > 0$: compute
            $b = at + v$, $\text{Area} = (ab + t)/2$, accumulate.
    \end{enumerate}
  \item Return the total sum.
\end{enumerate}

\section{Complexity}

\begin{itemize}
  \item \textbf{Orbit representative search:}
        $\displaystyle\sum_{m=1}^{\sqrt{n/2}} m \;\approx\; \frac{n}{4}$
        integer square root operations.
  \item \textbf{Orbit iteration:} $O(\log n)$ per orbit (exponential growth).
  \item \textbf{Total:} $O(n)$ integer square root operations.
  \item \textbf{Runtime:} $\sim$27\,s (C++), $\sim$5\,min (Python) for $n = 10^{10}$.
\end{itemize}

\section{Verification}

\begin{center}
\begin{tabular}{rr}
\toprule
$n$ & $S(n)$ \\
\midrule
50 & 9 \\
100 & 75 \\
500 & 576 \\
1\,000 & 1\,092 \\
5\,000 & 25\,845 \\
10\,000 & 39\,696 \\
\bottomrule
\end{tabular}
\end{center}

All values confirmed by independent brute-force computation.

\section{Answer}

\[
\boxed{2919133642971}
\]

\end{document}
